\section*{Подробные решения. Вариант 2}

\subsection*{1. Из шестнадцатеричной системы в восьмеричную}
Заменяем hex-цифры четвёрками битов и дополняем запись слева двумя нулями до групп по три:
\[
\begin{aligned}
F1C2_{16}&=\texttt{1111 0001 1100 0010}_2\\
&=\texttt{001 111 000 111 000 010}_2
=\boxed{170702_8}.
\end{aligned}
\]

\subsection*{2. Десятичное целое в шестнадцатеричной системе}
Делим на 16 с остатком:
\[
3178=16\cdot198+10,\qquad
198=16\cdot12+6,\qquad
12=16\cdot0+12.
\]
Читаем остатки снизу вверх: $12=C_{16}$, $10=A_{16}$. Ответ $\boxed{C6A_{16}}$. Проверка: $12\cdot256+6\cdot16+10=3178$.

\subsection*{3. Перевод десятичной дроби}
Умножаем остаток дробной части на 8, каждый раз записывая целую часть:
\[
\begin{array}{rcl@{\qquad}rcl}
0.667\cdot8&=&5.336&0.336\cdot8&=&2.688,\\
0.688\cdot8&=&5.504&0.504\cdot8&=&4.032.
\end{array}
\]
Первые четыре цифры --- \texttt{5254}. Следующая цифра получается из $0.032\cdot8=0.256$ и равна нулю, поэтому округление до четырёх знаков не меняет результат: $\boxed{0.5254_8}$.

\subsection*{4. Вычитание в шестёричной системе}
Выравниваем числа: $24453.12_6-00125.22_6$. В разряде шестых требуется заём из целой части: $1+6-2=5$. Целые части после займа: $24452_6-00125_6=24323_6$.
\[
\begin{array}{r@{}l}
24453&.12_6\\
-00125&.22_6\\\hline
24323&.50_6
\end{array}
\]
Ответ $\boxed{24323.50_6}$. Проверка дробной части: $0.12_6-0.22_6=-1/6$, а с занятым целым $1-1/6=0.50_6$.

\subsection*{5. Восстановление слагаемого $B$}
Дан код суммы: $177702_8=FFC2_{16}$. Старший бит равен 1, значит знаковое значение суммы --- $65474-65536=-62_{10}$. Вычитаем $A=-30$:
\[
B=-62-(-30)=-32_{10}=\boxed{-20_{16}}.
\]

Если под шестнадцатеричным значением подразумевается \emph{код} числа в памяти, он равен $2^{16}-32=\boxed{FFE0_{16}}$. Проверка: $-30\leftrightarrow FFE2_{16}$, и сумма кодов $FFE2_{16}+FFE0_{16}=1FFC2_{16}$; после отбрасывания 17-го бита получаем $FFC2_{16}$.

\subsection*{6. Вычитание в дополнительном коде}
$1011_8=512+8+1=521_{10}$, поэтому $A=-521$, $B=56$ и $A-B=-577$. В 16 битах коды $A=FDF7_{16}$, $-B=FFC8_{16}$:
\[
FDF7_{16}+FFC8_{16}=1FDBF_{16}
\ \longrightarrow\ \boxed{FDBF_{16}}.
\]
Перенос за пределы 16 битов отбрасывается. Проверка: $2^{16}-577=64959=FDBF_{16}$.

\subsection*{7. Одинарная точность IEEE 754}
Число положительно: знак $S=0$. Поскольку $512\le716.617<1024$, истинный порядок $E=9$; смещённый порядок $e=E+127=136=10001000_2$. Шаг формата при таком порядке равен $2^{9-23}=2^{-14}$. Для округления вычисляем
\[
716.617\cdot2^{14}=11741052+\frac{116}{125}.
\]
Поскольку $116/125>1/2$, берём ближайшее целое $11741053$. В 24-битном виде это \texttt{101100110010011101111101}. Старшую единицу не храним: поле дробной части --- \texttt{01100110010011101111101}. Соединяем поля:
\[
\texttt{0 10001000 01100110010011101111101}
=\boxed{4433277D_{16}}.
\]
При равенстве расстояний понадобилось бы правило чётности; здесь его применять не пришлось.

\subsection*{8. Функция логической схемы}
ИЛИ получает $A,B$, исключающее ИЛИ --- $A,C$, нижний НЕ --- $B$. Два элемента И и выходной НЕ дают
\[
\boxed{F=\lnot\bigl((A\lor B)\land(A\oplus C)\land\lnot B\bigr)}.
\]
Пересечения без точек не образуют дополнительных соединений.
